\documentclass[parskip=full-]{scrartcl}
\usepackage[utf8]{inputenc}
\usepackage[ngerman]{babel}
\usepackage{multicol}
\usepackage{amsmath}
\usepackage{amsfonts}
\usepackage{amssymb}
\usepackage{graphicx}
\usepackage{float}
\usepackage{enumitem} 
\usepackage{lmodern}
\usepackage{url}
\usepackage[pdftex,
    pdfauthor={Hoffmeier, Hoopmann, Johennecken},
    pdftitle={Anleitung Debian APT},
    pdfsubject={Anleitung Debian APT},
    pdfkeywords={Anleitung Debian APT},
    pdfproducer={Latex with hyperref},
pdfborder={0 0 0}]{hyperref}
\usepackage{caption}
\usepackage{subcaption}
\usepackage[skip=2pt,font=scriptsize]{caption}
\renewcommand*\familydefault{\sfdefault}


\renewcaptionname{ngerman}{\contentsname}{Inhalt}
\renewcaptionname{ngerman}{\listfigurename}{Abbildungen}
\renewcaptionname{ngerman}{\listtablename}{Tabellen}
\renewcaptionname{ngerman}{\figurename}{Abb.}
\renewcaptionname{ngerman}{\tablename}{Tab.}

\author{Hoffmeier, Hoopmann, Johennecken}
\title{Anleitung Debian APT}
\date {\today}

\usepackage{listings}
\usepackage{color}

\definecolor{mygreen}{rgb}{0,0.6,0}
\definecolor{mygray}{rgb}{0.5,0.5,0.5}
\definecolor{mymauve}{rgb}{0.58,0,0.82}

\lstset{ %
    backgroundcolor=\color{white},   % choose the background color; you must add \usepackage{color} or \usepackage{xcolor}
    basicstyle=\footnotesize,        % the size of the fonts that are used for the code
    breakatwhitespace=true,         % sets if automatic breaks should only happen at whitespace
    breaklines=true,                 % sets automatic line breaking
    captionpos=b,                    % sets the caption-position to bottom
    commentstyle=\color{mygreen},    % comment style
    deletekeywords={...},            % if you want to delete keywords from the given language
    %escapeinside={\%*}{*)},          % if you want to add LaTeX within your code
    extendedchars=true,              % lets you use non-ASCII characters; for 8-bits encodings only, does not work with UTF-8
    frame=single,                    % adds a frame around the code
    %keepspaces=true,                 % keeps spaces in text, useful for keeping indentation of code (possibly needs columns=flexible)
    keywordstyle=\color{blue},       % keyword style
    language=bash,                 % the language of the code
    otherkeywords={*,...},            % if you want to add more keywords to the set
    numbers=none,                    % where to put the line-numbers; possible values are (none, left, right)
    numbersep=5pt,                   % how far the line-numbers are from the code
    numberstyle=\tiny\color{mygray}, % the style that is used for the line-numbers
    rulecolor=\color{black},         % if not set, the frame-color may be changed on line-breaks within not-black text (e.g. comments (green here))
    showspaces=false,                % show spaces everywhere adding particular underscores; it overrides 'showstringspaces'
    showstringspaces=false,          % underline spaces within strings only
    showtabs=false,                  % show tabs within strings adding particular underscores
    stepnumber=1,                    % the step between two line-numbers. If it's 1, each line will be numbered
    stringstyle=\color{mymauve},     % string literal style
    tabsize=2,                       % sets default tabsize to 2 spaces
    title=\lstname                   % show the filename of files included with \lstinputlisting; also try caption instead of title
}


% Das Dokument beginnt hier
\begin{document}
\maketitle
\begin{figure}[H]
    \centering
    \includegraphics[scale=.5]{img/hs}
\end{figure}
\begin{abstract}
\end{abstract}
%	\tableofcontents
%\newpage
% Erstelle eine Hauptüberschrift
\section{VirtualBox Image}
Das Image für das Praktikum kann über TODO herunter geladen werden.

In VirtualBox den Punkt \glqq Appliance importieren\grqq aufrufen und Image über Assistenten einspielen.

Für das Praktikum wurde der Benuter \glqq john\grqq mit dem Password \glqq john\grqq angelegt.
\section{APT-Repoisitory hinzufügen}
Um das Repository von corsac (\url{http://molly.corsac.net/~corsac/debian/kernel-grsec/packages/}) in den Packetmanager zu integrieren muss die Datei \textit{/etc/apt/sources.list} verändert werden.

Dazu wird der folgende Text an die Datei an gehangen. 
\begin{lstlisting}
deb http://perso.corsac.net/~corsac/debian/kernel-grsec/packages/ jessie/
\end{lstlisting}

Zur Überprüfung der Dateien müssen der PGP Schlüssel von der Seite heruntergeladen und mit APT bekanntgemacht werden.
\begin{lstlisting}
wget -O - http://www.corsac.net/71ef0ba8.asc | apt-key -
\end{lstlisting}

Ein Update des APT-Baumes ermöglicht uns die Installation.
\begin{lstlisting}
apt-get update
\end{lstlisting}

Nun kann GrSec und PaX mit dem folgenden Kommando installiert werden.
\begin{lstlisting}
apt-get install linux-image-3.14.50-grsec+ paxctl paxtest
\end{lstlisting}

Der neue Kernel wird gleich in den Bootmanager grub eingetragen. Beim Neustart ist darauf zu achten, dass der GrSec-Kernel nicht der Standartkernel von grub ist. Da seine Versionsnummer niedriger als die des normalen Debian Kernels ist, muss dieser über die erweiterten Optionen gestartet werden.

Eine Übersicht über die Parameter (nur GrSec und PaX) die beim kompilieren des Kernels benutzt wurden, lässt sich mit dem folgenden Befehl ausgeben.
\begin{lstlisting}
cat /boot/config-3.14.50-grsec+ | grep -E '(PAX|GRKERN)'
\end{lstlisting}

Mit \glqq paxtest\grqq kann das PaX-System auf Funktion überprüft werden.
\begin{lstlisting}
paxtest kiddie
paxtest blackhat
\end{lstlisting}
\end{document}
